\centerline{\xiaoyi\bfseries Abstract}
\vspace{4em}

Federated learning enables distributed model training without moving raw
client data, but conventional digital aggregation requires communication
resources that grow with the number of participating clients. Analog
over-the-air federated learning (OTA-FL) addresses this bottleneck by
letting client updates superpose directly in the wireless multiple-access
channel. The same physical superposition that makes OTA-FL scalable also
injects channel noise and impulsive interference into the aggregated update.
This is especially challenging when the interference is heavy-tailed: rare
large perturbations can dominate a server update and destabilize plain
FedAvg-style training.

This thesis studies a PyTorch simulation framework for adaptive federated
learning over the air under symmetric alpha-stable interference. The AWGN
case is treated as the special case $\alpha=2$, while $\alpha<2$ models
heavier-tailed channel perturbations. The central hypothesis is that
server-side adaptive optimizers, especially AdaGrad-OTA and Adam-OTA, act as
implicit channel-noise filters: when the noisy aggregate has large
coordinate-wise magnitude, the fractional second-moment accumulator increases
and the effective learning rate automatically contracts.

The framework implements FedAvg-OTA, FedAvgM-OTA, AdaGrad-OTA, and Adam-OTA;
modular dataset partitioning; a noisy OTA aggregation oracle based on the
Chambers-Mallows-Stuck alpha-stable sampler; Median Anchored Clipping (MAC)
as robust pre-processing; JSON result export; and plotting scripts. The
revised experiments prioritize CIFAR-10 tail-index ablations and MAC
comparisons, with MNIST used as a lightweight sanity benchmark. The project
is not a full line-by-line reproduction of the reference paper, but it
implements the core heavy-tail setting needed to evaluate whether adaptive
server normalization improves OTA-FL robustness.

\noindent\textbf{Keywords:} federated learning; over-the-air aggregation;
alpha-stable noise; heavy-tailed interference; adaptive optimization; MAC
